\clearpage
\VerdierPhysicalPageAnchor{0206}%
\begingroup\small\setlength{\parskip}{0pt}\setlength{\emergencystretch}{1.5em}
\setlength{\abovedisplayskip}{4pt plus 2pt minus 2pt}%
\setlength{\belowdisplayskip}{4pt plus 2pt minus 2pt}%

\VerdierSemanticAnchor{ch3d-2-4-5}
\noindent\textbf{Remarque 2.4.5.}\quad
a) Soit $\mathcal A$ une catégorie abélienne dont tous les objets soient
isomorphes à des quotients d'objets projectifs. Un complexe de type
projectif n'est pas nécessairement $Q$-libre à gauche.

\VerdierSemanticAnchor{ch3d-2-4-5-example}
\emph{Exemple} : Soit $\mathcal A$ la catégorie des
$\mathbb Z/4\mathbb Z$-modules. Le complexe :
\[
\cdots\longrightarrow\mathbb Z/4\mathbb Z
\longrightarrow\mathbb Z/4\mathbb Z
\longrightarrow\mathbb Z/4\mathbb Z
\longrightarrow\mathbb Z/4\mathbb Z
\longrightarrow\mathbb Z/4\mathbb Z\longrightarrow\cdots
\]
dont toutes les différentielles sont la multiplication par $2$ est
acyclique et de type projectif. Il n'est cependant pas homotope à zéro. Le
même complexe fournit un exemple de complexe acyclique de type injectif,
non homotope à zéro.

\VerdierSemanticAnchor{ch3d-2-4-5-b}
b) Même lorsque la dimension projective des objets de $\mathcal A$ n'est
pas bornée, il peut exister des complexes de type projectif $Q$-libres à
gauche dont l'image dans $D(\mathcal A)$ n'appartienne pas à
$D^{-}(\mathcal A)$. Par exemple, un complexe projectif de
Cartan-Eilenberg est un objet $Q$-libre à gauche. Plus généralement, toute
somme directe d'objets $Q$-libres à gauche est $Q$-libre à gauche.

\VerdierSemanticAnchor{ch3d-2-4-6}
\medskip\noindent\textbf{2.4.6.}\quad
Nous laissons au lecteur le soin de formuler les énoncés qu'on obtient en
appliquant les résultats de ce numéro à la catégorie opposée à la catégorie
$\mathcal A$.
\endgroup

